La lettre reproduite ci-dessous présente les objectifs généraux de la
démarche ainsi que les conditions de sa mise en œuvre.

\begin{tcolorbox}[
breakable,
colback=gray!5,
colframe=black,
boxrule=0.5pt,
arc=3mm
]
\setlength{\parskip}{0.5em}
IHEDN AR 16

Comité d'études « Crises majeures »

\hfill Paris, le 6 mars 2026

Chère Présidente,\\
Chères et Chers Membres CoDir,

Le 50e anniversaire de notre Association constitue un moment important de son histoire. Au-delà de la célébration légitime du parcours accompli, cet anniversaire peut aussi devenir une occasion de projection vers l'avenir.

Dans un contexte marqué par la multiplication des crises, les réseaux d'expertise et de confiance constituent une ressource précieuse pour la société. L'AR 16 représente, à cet égard, un capital humain et relationnel exceptionnel, encore insuffisamment cartographié.

Notre Comité d'études pose la question : que peut apporter notre réseau d'auditeurs en situation de crise majeure ?

Pour répondre à cette interrogation, nous proposons le lancement, à l'occasion du 50e anniversaire, d'une initiative structurante :

\begin{itemize}
    \item la cartographie des compétences, disponibilités et engagements des membres de l'AR 16.
\end{itemize}

Cette démarche prendra la forme d'une enquête auprès des membres, permettant de :

\begin{itemize}
    \item mieux connaître les profils et expertises des auditeurs ;
    \item identifier la disponibilité et le mode d'engagement en cas de crise ;
    \item structurer des groupes de réflexion ou d'analyse ad hoc.
\end{itemize}

Ainsi, notre Association pourrait contribuer à renforcer la culture de coopération et de résilience, sans jamais se substituer aux responsabilités des acteurs publics.

L'enquête - sous la forme d'un questionnaire anonyme - devient un instrument de connaissance du capital stratégique de l'AR 16. Elle permet aussi de transformer l'expérience individuelle en connaissance collective.

Il ne s'agit ni d'un contrôle, ni d'une obligation, ni d'une évaluation. Les réponses seront traitées de manière anonyme, confidentielle et analysées globalement.

Pour nous accompagner dans cette démarche, le Comité d'études a échangé avec le Dr Victor Violier, sociologue à l'\gls{irsem}. La faisabilité du projet suppose toutefois le concours d'un délégué à la protection des données afin de garantir la conformité au règlement européen sur la protection des données.

Dans cette perspective, un rapprochement avec l'Union IHEDN, dépositaire des fichiers des associations régionales, pourrait être envisagé.

L'enquête, dont le lancement pourrait intervenir dans le courant du premier semestre 2026, ne donnera lieu à une exploitation complète des résultats qu'au cours de l'année suivante. Elle pourrait néanmoins s'inscrire dans le cadre du cinquantième anniversaire de l'AR~16 et faire l'objet d'une restitution publique, éventuellement avec le concours de l'\gls{irsem}, lors d'un événement organisé en 2027.

Notre ambition est réelle : l'enquête permet de célébrer notre histoire, mais aussi, dans le cadre d'une démarche participative, de réfléchir collectivement à notre avenir.

Nous vous remercions, Chère Présidente, Chères et Chers Membres du CoDir, de votre accueil, compréhension et confiance.

\begin{flushright}
Les Membres du Comité « Crises majeures »

(Présents à la réunion du 5 mars 2026)

Valérie Arlé-Faivre, Gérard Turck, Aurélien Duvignac-Rosa,\\
Jean-Jacques Cleynen, Bruno Leray, Alain Fontan
\end{flushright}

\end{tcolorbox}